% Proposed Framework: AnimatorSBT

\textbf{AnimatorSBT} integrates Soulbound Token issuance into the anime production
workflow through a production management tool the animator already uses.
% ----- 4.1 System Architecture -----
\subsection{System Architecture}
\label{subsec:architecture}

Four components carry the design. A production management client---the
application an animator already uses to track cut assignments, status and
deadlines---is the source of the contribution record. An issuance service
aggregates those entries at a milestone, pins the metadata, and calls the
contract. The contract itself is an ERC-721 token~\cite{entriken2018erc721} on
Polygon PoS with transfers disabled, exposing minting, revocation and metadata
retrieval. A public verification portal resolves an address to the certificates it
holds. Only the first of these is visible to the animator; the rest run behind it,
which is a usability decision with a disclosure obligation attached
(\cref{subsec:adoption_barriers}).

The components are not what needs defending. Three choices are.

\textbf{Why blockchain rather than a centralized database?}
The conventional answer is a database run by one party, a studio or us. It fails the
requirements of \cref{subsec:formalization} in three places, and the chain rescues
each only partly.

A database goes with its operator, and studios close. What survives on chain is
ownership, the issuance timestamp and the metadata pointer, whether or not we do. What
does not is the human-readable credential, because the metadata sits off chain and
depends on someone continuing to pin it, and because the portal and issuance service
are ours to run. That a certificate exists and whose it is outlives the issuer; its
content does not, unless another party keeps it.

An animator works across dozens of studios, each holding a fragment in its own system
with no way to combine them. Certificates accumulate instead at one address, so a
career is presentable as a link rather than assembled from testimony, subject to the
authenticity caveat throughout this paper.

An administrator can rewrite or delete records; an on-chain record cannot be altered
afterwards without leaving a trace. This is narrower than trustless, since issuance is
permissioned and only role-holders can mint. The animator gets not a system nobody
controls, but one in which what was issued, when and to whom cannot be changed
quietly.

\textbf{Why SBTs rather than Verifiable Credentials?}
W3C Verifiable Credentials~\cite{sporny2022vc}, the credential layer of the
self-sovereign identity stack M{\"u}hle et al.\ survey~\cite{muhle2018ssi}, are the
obvious alternative: mature, capable of selective disclosure, stored off chain, and
suited to a durable issuer presenting credentials bilaterally. They are not
unworkable here, and a VC anchored on chain by a hash or a status entry would not
depend on the issuer still answering. Our reason for choosing SBTs is ergonomic and
concerns R4 and R5. Certificates accumulate at one address and are queryable in one
call, where the VC model asks the holder to store, manage and present each credential
and routes verification through the issuer's key or a status registry the issuer
maintains. For a holder-passive experience aimed at non-technical animators, the
address is the portfolio.

Kang and Lemieux build the fuller version of that stack, adding encryption and usage
control so that a data subject licenses access and the licence itself constrains
secondary use~\cite{kang2021ssi}. The contrast is instructive: theirs governs what may
be done with data a subject controls, ours is a narrow attestation the animator wants
read as widely as possible, and usage control is the wrong tool for a credential whose
value is its publicity. The two families are not exclusive. A VC wrapper over the same
metadata would add interoperability, and per-attribute disclosure of the kind the VC
ecosystem supports would be finer than the commitment of \cref{subsec:privacy}, which
reveals the name in full or not at all.

\textbf{Why a bespoke contract rather than a general attestation service or a
signed log?}
A generic on-chain attestation registry such as the Ethereum Attestation
Service~\cite{eas} could carry these attestations without a custom contract, and is a
viable substrate: the co-signing design of \cref{subsec:cosigning} could be expressed
as EAS attestations. We chose a dedicated ERC-721 contract because it bundles
soulbound semantics, role-gated issuance and role-gated co-signing into one auditable
unit depending on no external registry's schema or continued operation, and because a
non-transferable token gives the portfolio-as-balance ergonomic natively. Those are
arguments from self-containment, not from technical superiority.

Wang et al.\ question the binary itself, proposing bounded rather than prohibited
transfer as a limited transfer count~\cite{wang2026counted}. For a credential of this
kind the binary is right: a work record that can change hands even a fixed number of
times records whose account holds it, not who did the work, and what the animator
needs is precisely that it cannot be sold when they are short of money.

One could also forgo the chain and publish a signed append-only log, as certificate
transparency does. That achieves tamper-evidence more cheaply and fails on the
requirement that matters here: verification depends on the operator continuing to
serve the log and its key, where a public-chain record is checkable through any node
whether or not the issuer exists. For integrity alone the signed log is the better
answer; for persistence beyond the issuer it is not an answer at all.

% ----- 4.2 SBT Metadata Schema -----
\subsection{SBT Metadata Schema}
\label{subsec:metadata}

Each certificate records one project contribution. The metadata is a JSON document
on IPFS~\cite{benet2014ipfs}, which addresses content by a hash of itself rather than
by location, and the on-chain token URI carries that hash. \Cref{tab:metadata} defines
the schema.

\begin{table}[ht]
\centering
\caption{AnimatorSBT metadata schema. The JSON follows an NFT-compatible
structure with top-level \texttt{name} and \texttt{description} fields;
contribution data is stored under an \texttt{attributes} object.}
\label{tab:metadata}
\footnotesize
\begin{tabular}{@{}llp{3.0cm}@{}}
\toprule
\textbf{Field} & \textbf{Type} & \textbf{Description} \\
\midrule
\texttt{name}              & string   & Display name for the SBT \\
\texttt{description}       & string   & Human-readable summary \\
\midrule
\multicolumn{3}{@{}l}{\textit{attributes} (contribution data):} \\
\texttt{animator\_name}    & string   & Held by the holder, not published (\cref{subsec:privacy}) \\
\texttt{animator\_wallet}  & address  & Animator's wallet address \\
\texttt{project}           & string   & Anime project or episode title \\
\texttt{role}              & enum     & e.g., \texttt{key\_anim}, \texttt{in\_between}, \texttt{coloring} \\
\texttt{task\_count}       & integer  & Tasks completed (cuts, frames) \\
\texttt{start\_date}       & date     & Contribution start date \\
\texttt{end\_date}         & date     & Contribution end date \\
\texttt{studio\_name}      & string   & Contracting studio (nullable) \\
\texttt{issued\_at}        & datetime & ISO~8601 issuance timestamp \\
\midrule
\texttt{evidence\_files}   & array    & IPFS URIs of supporting files \\
\texttt{salt}              & bytes32  & Held by the holder, not published \\
\texttt{evidence\_hash}    & bytes32  & Keccak-256 over the canonical private record \\
\bottomrule
\end{tabular}
\end{table}

Storing the record on chain was never an option at this size. Sward et al.\ survey
the alternatives for binding off-chain data to a ledger, with their cost and permanence
trade-offs~\cite{sward2018insertion}; we take the cheapest, anchoring a hash and
keeping the content elsewhere, and \cref{sec:evaluation} shows where the cost lands
instead.

Two fields carry the design rather than the data. \texttt{evidence\_hash} links the
certificate to the work log it was built from, so a verifier can check correspondence
without being shown the log; \cref{subsec:verification} sets out how little that
establishes. \texttt{studio\_name} is nullable, and the consequence is sharper than it
looks: a certificate with no named counterparty has nobody who could co-sign it, so the
privacy affordance is available exactly on the confidential, high-value productions
where third-party verification would be worth most.

% ----- 4.3 Issuance Flow -----
\subsection{Issuance Flow}
\label{subsec:flow}

Issuance follows the animator's ordinary work. They record cut assignments, status
changes and completion times in the management client as they already do. At a
milestone---delivery of an episode, or a request from the animator---the client asks
the issuance service to certify the period. The service aggregates the relevant log
entries, computes the evidence hash over their canonical serialisation
(\cref{subsec:canonical}), pins the metadata, and calls \texttt{mint()} with the
animator's address and the resulting URI. A certificate appears. In the production
design a meta-transaction relayer sponsors gas, so the animator neither holds tokens
nor signs a transaction; the prototype pays gas from an operational wallet.

The resulting record is narrower than it looks. Suppose a key animator is assigned
25 cuts for an episode through a secondary subcontractor, works ten weeks, and is
never named in the broadcast credits. The certificate establishes that they
declared 25 key-animation cuts on a production in that period, and that the
declaration existed at that time and has not changed since. It does not establish
that they did the work. That distinction is the whole subject of
\cref{subsec:cosigning}: with the subcontractor's co-signature the same record
becomes an attestation by the party that accepted the work, and without it the
certificate is a timestamped self-report---better than an unverifiable line on a
r\'esum\'e, and less than a credential.

% ----- 4.4 Verification Mechanism -----
\subsection{Verification Mechanism}
\label{subsec:verification}

A verifier with an animator's address queries the contract for the tokens it holds
and their metadata URIs. Because the tokens cannot be transferred, holding one is
evidence that the holder is who it was issued to. A portal renders the resulting
portfolio and can recompute the evidence hash, with the animator's consent,
against the work log the certificate was built from.

What that recomputation establishes is narrow. From the pinned metadata alone the
hash detects whether the file changed after minting and nothing more, because the
issuer authored both the attributes and the hash over them. Against the animator's own
local log it is stronger, tying the certificate to a contemporaneous record, but still
integrity rather than authenticity: the log is self-reported, so a match shows the
certificate agrees with what was claimed, not with what was done. It is also fragile,
requiring consent and the survival of local data on a device. Authenticity is the
co-signature's job (\cref{subsec:cosigning}), and anchoring on an accepted order
(\cref{subsec:integration}) is what gives the co-signature something already
determined to sign.

% ----- 4.5 Privacy Considerations -----
\subsection{Privacy Considerations}
\label{subsec:privacy}

Privacy here turns on a question the schema got wrong: what has to be published for
a certificate to be useful.

The prototype stored \texttt{animator\_name} in cleartext on public IPFS, which was
indefensible and is no longer the default. A revocation flag does not undo writing a
name to permanent, content-addressed storage that an on-chain record points at, and
replacing the name with a bare hash does not help either. The Article~29 Working Party
treats hashing as pseudonymisation rather than anonymisation, and notes that a known
range of inputs can be replayed through the function to find which one produced a given
digest~\cite{art29wp2014anonymisation}; Finck applies that reasoning to ledgers,
concluding that hashed data on a chain remains personal
data~\cite{finck2018blockchains,finck2019stoa}. The set of names in one industry is
small enough for the replay to be practical.

The certificate now commits to a record that names the holder and publishes the
record without the name. The evidence hash is taken over the full private
record---the public attributes, the holder's name, and a random salt the holder
keeps---while what is pinned carries only role, volume, period and the address.
A verifier shown the name and the salt recomputes the hash and learns that this
certificate was issued for that person. Anyone else sees a certified role and
volume bound to an address and cannot search for the name, because they do not
have the salt. The salt is what closes the replay described above: the candidates an
attacker must test are no longer the names in an industry but the values of a
256-bit salt, which is the distinction the Working Party draws when it separates a
plain hash from a keyed one whose key is unavailable~\cite{art29wp2014anonymisation}.

Linkability then comes from one place instead of two. The certificates a person
holds are linkable because they sit at one address, which is what makes a
portfolio a portfolio; distinct salts mean the digests add nothing to that. The
holder can therefore decide what an address is worth revealing about them, rather
than having each certificate publish the answer independently.

This is not zero knowledge: the holder reveals the name in full to whoever they
choose to show it to. What changes is who discloses, and how often. Disclosure becomes
an act the holder repeats per verifier instead of a publication the issuer performs
once, permanently, for everyone. Ten tests cover the construction, including that a
substituted name and that altered public attributes both fail against a valid
certificate.

What this does not fix is inference. Role, volume, period and a co-signing studio may
identify someone in a small industry whether or not the name is published, so we do not
claim the published record is anonymous. The industry's own convention points the same
way: the state's model contract provides for crediting under a stage or pen name rather
than a legal one~\cite{bunka2024model}, so a certificate carrying a professional
identity behaves like a credit rather than departing from one.

Here we part company with the recent literature. Lyu et al.\ bind a soulbound token to
an identity engine spawning per-service child identities, precisely so that verifiers
cannot link a holder's credentials to one another~\cite{lyu2026ssi}. That is the right
goal for credentials presented bilaterally and the wrong one here. A portfolio is
valuable because its entries gather at one identity, and severing that link would remove
the property the animator needs, which is why we hide the name and keep the address
rather than the reverse. Zero-knowledge techniques belong in this design at the level of
attributes rather than holders---proving a volume of work across a number of productions
without naming them (\cref{subsec:trust})---not in severing the link that makes a career
record a record.

Three smaller measures apply beyond the naming question. Only the token identifier,
owner address and metadata URI go on chain. Fields covered by a non-disclosure
agreement, such as the studio and the production title, can be omitted while still
certifying role and volume, at the cost discussed in \cref{subsec:adoption_barriers}.
And account abstraction keeps animators from handling private keys.

Erasure we cannot engineer around. \texttt{revoke()} sets a flag rather than burning
the token or clearing its URI, because preserving the audit trail is the point of
putting the record on a chain at all, so a revoked token stays queryable and its URI
still resolves. Erasing the content means unpinning the metadata, which makes it
unreachable through public gateways but not gone: anyone who kept a copy can re-pin it,
and the token, its timestamp and the revocation record persist regardless. Encrypting
the metadata and discarding the key would give a genuine erasure path, and we did not
implement it. The legal literature reaches the same wall. Finck's study for the European
Parliament puts the append-only ledger against Articles~16 and~17 of the GDPR and finds
the tension fundamental rather than incidental, since the property that makes
modification onerous is the property the design is chosen for~\cite{finck2019stoa}. We
therefore claim a partial, best-effort capability rather than a right to be forgotten,
and treat the tension with the Act on the Protection of Personal Information as
unresolved (\cref{subsec:legal}).

% ----- 4.6 Co-signing / Attestation Design -----
\subsection{Co-signing and the Authenticity Layer}
\label{subsec:cosigning}

Everything above produces a record that is tamper-evident, timestamped and
portable, and none of it makes the record true. This subsection designs the layer
that would, and implements it far enough to measure.

\paragraph{What a credential is worth depends on who signed it.}
A self-attested certificate, minted from the animator's own log, establishes existence,
integrity and a timestamp: more than an unverifiable claim, less than an authenticated
one. A studio-attested certificate, countersigned by the party that accepted the work,
is the level a hiring verifier can weigh. A middle level is conceivable---co-workers on
the same production endorsing a claim, a social signal short of an institutional
one---but we have not built it, so the contract's \texttt{attestationLevel} is binary
rather than graded. A verifier should read the level, not just the token.

\paragraph{Reference implementation.}
This reuses the contract's role-based access control (\cref{subsec:contract}). \texttt{AnimatorSBTV2} adds a
\texttt{STUDIO\_ATTESTER\_ROLE} and the function
\texttt{coSign(uint256 tokenId)}, which records an attestation by
\texttt{msg.sender} against an existing token and emits \texttt{SBTCoSigned}.
We bind the attester to \texttt{msg.sender} rather than passing it as an argument
(as an earlier sketch did) so that the on-chain signature itself proves who
attested, with no spoofable parameter. The mechanism is idempotent (the same
attester cannot co-sign twice), supports multiple distinct studios co-signing one
token, and refuses to co-sign a revoked token. Studios are granted the attester role
through the same admin pathway used for issuers, and \cref{subsec:canonical} reports
what the contract costs to deploy and to co-sign.

\paragraph{The attester-independence caveat.}
The authenticity gain from co-signing is exactly as strong as the attester's
independence and real-world identity binding---neither of which the contract can
enforce on its own. Two gaps remain open. First, in the current single-operator MVP
the issuer wallet also holds \texttt{STUDIO\_ATTESTER\_ROLE} (and the issuance
service defaults the attester key to the issuer key), so a ``studio-attested'' token
can in fact be self-co-signed by the operator; this configuration adds no
authenticity and exists only to exercise the mechanism end to end. Second, even with
a distinct attester wallet, nothing \emph{on-chain} binds \texttt{STUDIO\_ATTESTER\_ROLE}
to a verified legal studio: the admin (TOONIQ) decides who is treated as a studio, so
co-signing shifts the root of trust from the animator's self-report to TOONIQ's
off-chain vetting rather than removing trust. The second gap has been attacked
directly: Lauinger et al.\ let a holder verify that a credential issuer was itself
legitimately authorised, using an accumulator and a zero-knowledge proof rather than a
named administrator~\cite{lauinger2021apoa}. That is the cryptographic form of what we
do administratively and the direction a deployment ought to take; we do not adopt it
because our attester set is small enough to name, and because it would add machinery to
a mechanism we have not operated at all. Closing both gaps---an independent attester
operated by a real studio, with an auditable identity binding---is what would convert a
co-signed record from \emph{operator-attested} to \emph{independently authenticated}
(\cref{sec:conclusion}). Until then what this paper contributes is a tamper-evident,
portable contribution ledger the holder controls, together with a tested reference
implementation of the authenticity layer whose independent operation would make it a
credential.
